%! Multiplying \& Dividing Rational Expressions — Unit 4, Lesson 4.2 — Homework (Answer Key)
\documentclass[10pt]{article}
\usepackage{algebra2-article}
\usepackage{algebra2-key}   % pulls in -boxes; do NOT also load -boxes
\newcommand{\TallMath}[1]{$\displaystyle #1\rule[-1.4em]{0pt}{3.2em}$}

\begin{document}

\pageheader{Unit 4, Lesson 4.2}{Homework --- Answer Key}
\namedateperiod

\vspace{0.10in}

\begin{notesbox}{Practice}
Every answer needs \emph{two} parts: the simplest form \textbf{and} the restrictions. Check all
\textbf{three} sources on every division --- both denominators \emph{and} the divisor's numerator.

\begin{enumerate}[leftmargin=1.9em, itemsep=11pt]
  \item \textbf{Monomial factors.} Simplify and restrict.
    \begin{enumerate}[label=(\alph*), leftmargin=2.1em, itemsep=5pt]
      \item $\dfrac{4x^3}{9y}\cdot\dfrac{3y^4}{8x}=\;$\ans{$\frac{x^2y^3}{6}$} \quad restrictions: \ans{$x\neq0,\ y\neq0$}
      \item $\dfrac{5a^2}{6b^3}\div\dfrac{10a}{3b}=\;$\ans{$\frac{a}{4b^2}$} \quad restrictions: \ans{$a\neq0,\ b\neq0$}
      \item $\dfrac{7x^2y}{2}\div\dfrac{14xy^3}{5}=\;$\ans{$\frac{5x}{4y^2}$} \quad restrictions: \ans{$x\neq0,\ y\neq0$}
    \end{enumerate}
    \vspace{2pt}
    In part (c), neither original denominator can ever be zero. So where did the restrictions come
    from? \ans{the divisor's numerator $14xy^3$ --- source 3}

  \item \textbf{Binomial and quadratic factors.} Factor completely first. One of these hides a pair of
        \emph{opposite} binomials --- watch the sign.
    \begin{enumerate}[label=(\alph*), leftmargin=2.1em, itemsep=7pt]
      \item $\dfrac{x+3}{x-2}\cdot\dfrac{x^2-4}{x^2-9}=\;$\ans{$\frac{x+2}{x-3}$} \quad $x\neq$ \ans{$2,\ 3,\ -3$}
      \item $\dfrac{x^2+6x+8}{x^2-1}\cdot\dfrac{x-1}{x+4}=\;$\ans{$\frac{x+2}{x+1}$} \quad $x\neq$ \ans{$1,\ -1,\ -4$}
      \item $\dfrac{x^2-36}{5x}\div\dfrac{x+6}{10x^2}=\;$\ans{$2x(x-6)$} \quad $x\neq$ \ans{$0,\ -6$}
      \item $\dfrac{x^2+2x-15}{x^2+x-6}\div\dfrac{x+5}{x+3}=\;$\ans{$\frac{x-3}{x-2}$} \quad $x\neq$ \ans{$-3,\ 2,\ -5$}
      \item $\dfrac{x^2-49}{x+1}\cdot\dfrac{3x+3}{7-x}=\;$\ans{$-3(x+7)$} \quad $x\neq$ \ans{$-1,\ 7$}
      \item $\dfrac{x^2-2x}{x^2+5x+6}\div\dfrac{x^2-4}{x+3}=\;$\ans{$\frac{x}{(x+2)^2}$} \quad $x\neq$ \ans{$-2,\ -3,\ 2$}
    \end{enumerate}

  \item \textbf{A rectangle.} A rectangle has width $\dfrac{2x}{x+5}$ units and length
        $\dfrac{x^2+5x}{4}$ units.
    \begin{enumerate}[label=(\alph*), leftmargin=2.1em, itemsep=5pt]
      \item Area $=$ width $\times$ length $=$ \ans{$\frac{x^2}{2}$}, \; with restriction $x\neq$ \ans{$-5$}
      \item Your area has no denominator containing $x$. Explain why the restriction still belongs to
            it. \ans{at $x=-5$ the width itself has no value, so there is no rectangle to measure}
      \item Both dimensions must be positive for a real rectangle. Which values of $x$ make sense?
            Write your answer as a union of intervals. \ans{$(-\infty,-5)\cup(0,\infty)$}
    \end{enumerate}

  \item \textbf{Justify.} Are $\dfrac{x+1}{x-2}\div\dfrac{x+1}{x-3}$ and $\dfrac{x-3}{x-2}$ the same
        expression? Answer in terms of \emph{inputs}: where do they agree, where do they disagree, and
        what does each graph look like?
        \ansline{They are \emph{equivalent} --- keep--change--flip gives $\frac{x+1}{x-2}\cdot\frac{x-3}{x+1}=\frac{x-3}{x-2}$ --- and they agree at every}
        \ansline{input both accept. But the quotient also excludes $x=3$ (the divisor's denominator) and $x=-1$ (the divisor's}
        \ansline{numerator), so its graph is the graph of $\frac{x-3}{x-2}$ with two extra holes punched out, at $(3,0)$ and $\left(-1,\frac43\right)$.}
\end{enumerate}
\end{notesbox}

\begin{extensionbox}
\textbf{Part 1 --- Three students, three answers.} Each student simplified
$\dfrac{x^2-1}{x+4}\div\dfrac{x-1}{x^2-16}$ and got something different. Decide who is right, and name
each error.
\begin{center}
\renewcommand{\arraystretch}{1.4}
\begin{tabularx}{\linewidth}{@{}l l X@{}}
\hline
\textbf{Student} & \textbf{Answer} & \textbf{Right or wrong --- and why?} \\
\hline
Dev   & $(x+1)(x-4)$, $x\neq-4,4,1$ & \ans{Correct --- all three sources found.} \\
Elin  & $(x+1)(x-4)$, no restrictions & \ans{Incomplete: right form, but a polynomial answer still carries all three.} \\
Farid & $\dfrac{(x-1)^2(x+1)}{(x+4)^2(x-4)}$ & \ans{Wrong: multiplied by the divisor instead of by its reciprocal --- never flipped.} \\
\hline
\end{tabularx}
\end{center}

\vspace{4pt}
\textbf{Part 2 --- Design one.} Write a \emph{division} problem whose simplified answer is $x+2$ and
whose restrictions are exactly $x\neq3$, $x\neq0$, and $x\neq-1$.
\begin{center}
\ans{$\frac{x(x+1)(x+2)}{x-3}$} $\div$ \ans{$\frac{x(x+1)}{x-3}$}
\end{center}
\begin{enumerate}[label=(\alph*), leftmargin=1.9em, itemsep=5pt]
  \item Which factor did you put in the divisor's \emph{numerator}, and which restriction did that buy
        you?
  \item Your answer $x+2$ is a polynomial --- its graph is a line. What does that line look like once
        the restrictions are applied?
\end{enumerate}
\ansline{(a) $x(x+1)$ --- on top of the divisor it forces $x\neq0$ and $x\neq-1$, and after the flip it lands downstairs}
\ansline{where the matching factors upstairs divide it out. The $(x-3)$ in both denominators supplies $x\neq3$.}
\ansline{(b) the line $y=x+2$ with three holes punched out, at $(3,5)$, $(0,2)$, and $(-1,1)$.}
\end{extensionbox}

\begin{spiralbox}
\textbf{Coming next --- Lesson 4.3: Adding \& Subtracting Rational Expressions (and complex
fractions).} Multiplying and dividing never needed a common denominator; adding and subtracting need
nothing else. You will factor every denominator (same first step as always), build the \textbf{least
common denominator} out of those factors, and rewrite each fraction to match it. Then the payoff:
a \textbf{complex fraction} --- a fraction whose numerator and denominator are themselves fractions ---
is nothing more than \emph{combine the top, combine the bottom, then divide}. And dividing is exactly
what you did today, three restriction sources and all.
\end{spiralbox}

\begin{teachernote}
\textbf{Item 1(c)} is the item to debrief first: the two written denominators are the constants $2$ and
$5$, so \emph{every} restriction on that problem comes from the divisor's numerator. A student who
answers ``none'' has today's whole idea still missing.
\textbf{Item 2(e)} is the Lesson 4.1 sign check --- $7-x=-(x-7)$, so the answer is $-3(x+7)$; a missing
$-1$ is the most common single lost point on this set. \textbf{Item 2(f)} is the hardest: the divisor's
numerator $x^2-4$ contributes \emph{two} values ($x=2$ and $x=-2$), one of which duplicates a
denominator restriction, and the answer keeps a squared factor $(x+2)^2$.

\vspace{3pt}
\textbf{Item 3} is worth debrief time. The area simplifies to $\frac{x^2}{2}$ --- an expression with no
$x$ in any denominator --- and students still have to carry $x\neq-5$. Part (c) surprises them: both
dimensions are also positive for $x<-5$ (negative over negative), so the feasible set is a union, not a
ray. Accept $(0,\infty)$ with sound reasoning, but surface the second interval.

\vspace{3pt}
\textbf{Extension Part 2} accepts many answers; check three things --- the answer really does simplify
to $x+2$, the divisor's numerator supplies $x\neq0$ and $x\neq-1$, and every factor placed downstairs
has a partner upstairs so nothing is left over. Part (b) previews 4.5 directly: a polynomial answer
with three holes is a graph students will be asked to produce from an equation.
\end{teachernote}

\end{document}
